\documentclass[11pt,a4paper]{article}

% ── Packages ──
\usepackage[utf8]{inputenc}
\usepackage[T1]{fontenc}
\usepackage{booktabs}
\usepackage{amsmath,amssymb}
\usepackage{graphicx}
\usepackage[margin=2.5cm]{geometry}
\usepackage{caption}
\usepackage{hyperref}
\usepackage{multirow}
\usepackage{array}
\usepackage{pifont}

% ── Title ──
\title{ProCyon v2: Learning Robust Microbiome Representations for\\Accurate IBD Prediction and Evidence-Grounded Biological Interpretation}
\author{}
\date{}

\begin{document}
\maketitle

\begin{abstract}
We present ProCyon v2, a microbiome analysis framework that decouples disease classification from biological interpretation. Through systematic architecture ablation on the clean\_2538 IBD dataset (659 train, 167 test), we find that: (1) a minimal encoder—nn.Embedding with masked mean pooling, no pretraining, no Transformer—achieves 92.57\% IBD diagnosis accuracy (AUC=0.975), significantly outperforming a 34M-parameter pretrained MGM Transformer encoder by 41.7 percentage points (92.57\% vs 50.9\%); (2) XGBoost on raw genus abundance is a strong baseline (92.22\%), but lacks the embedding structure needed for interpretability; (3) embedding dimension saturates at $E{=}512$ (93.41\%), requiring only 760K parameters; (4) cross-dataset transfer evaluation shows training on $5\times$ more diverse data improves generalization from 61.8\% to 88.6\% ACC; (5) LOO Spearman rank correlation $\rho{=}0.72$ across folds confirms consistent genus importance ranking, with deletion tests showing 5.6$\times$ larger AUC drop than random; (6) Qwen2-7B-Instruct fails as a classifier (70.7\% $\pm$ 11.9\% instability) but excels as an interpreter when grounded by LOO attributions—halving hallucination (0.72 to 0.36 genera/response), achieving 98\% prediction consistency and 0.94 direction accuracy on all 167 test samples. ProCyon v2 establishes that the model discovers patterns, LOO attribution provides evidence, and the LLM explains the evidence.
\end{abstract}

% ═══════════════════════════════════════
% 1. INTRODUCTION
% ═══════════════════════════════════════
\section{Introduction}

Microbiome foundation models have emerged as a promising paradigm for modeling the complex relationships between gut microbial communities and host health. Recent models including MGM~\cite{mgm2024}, Waypoint~\cite{waypoint2026}, and ProCyon~\cite{procyon2025} follow a common blueprint: tokenize genus abundance profiles, encode them with a Transformer, and fine-tune on downstream tasks.

However, three critical questions remain under-explored:

\begin{enumerate}
  \item \textbf{Architecture necessity}: Do genus-level microbiome classification tasks actually benefit from Transformer encoders and large-scale pretraining?
  \item \textbf{LLM role}: Should large language models serve as classifiers, or as reasoning interfaces on top of a specialized classifier?
  \item \textbf{Evidence grounding}: Can model explanations be verified against literature and made robust across data splits?
\end{enumerate}

This paper systematically investigates these questions through controlled experiments. Our findings challenge several assumptions in current microbiome foundation model design and lead to ProCyon v2: a simple embedding-based classifier with LOO-attributed feature importance, where the LLM serves solely as an interpreter—not a predictor.

\textbf{Contributions:}
\begin{itemize}
  \item \textbf{Architecture minimalism}: SimpleEmb (nn.Embedding + mean pool + MLP, 1.1M params, random init) achieves 92.57\% ACC, exceeding MGM's 34M-param pretrained Transformer (50.9\%) by 41.7pp
  \item \textbf{Comprehensive baseline comparison}: 9 methods compared on identical train/test split with 6 performance metrics
  \item \textbf{Component ablation}: Isolate contributions of embedding, random projection, non-linearity, and end-to-end training
  \item \textbf{Embedding dimension analysis}: Performance saturates at $E{=}512$ (93.41\% ACC, 760K params)
  \item \textbf{Cross-dataset transfer}: Training on $5\times$ larger diverse data improves cross-dataset generalization by 26.8pp
  \item \textbf{LOO Attribution Reliability}: Spearman $\rho{=}0.72$ rank consistency, deletion test validation, literature direction agreement (5/6)
  \item \textbf{LLM explanation benchmark} ($n{=}167$): LOO grounding halves hallucination (0.72 to 0.36/response), achieves 98\% consistency and 0.94 direction accuracy
\end{itemize}

% ═══════════════════════════════════════
% 2. RELATED WORK
% ═══════════════════════════════════════
\section{Related Work}

\textbf{MGM}~\cite{mgm2024} introduced the first large-scale microbiome foundation model: an 8-layer Transformer pretrained with autoregressive next-genus prediction on MicroCorpus-260K (263,302 samples, 9,665 genera). MGM demonstrated cross-regional diagnosis, keystone genus discovery via attention analysis, and in silico perturbation experiments. Its rank-based encoding mitigates extreme abundance values but discards relative abundance information.

\textbf{BiomeGPT}~\cite{biomegpt2026} is a transformer-based foundation model pretrained on 13,300 species-level human gut metagenomes spanning 32 phenotypes. It uses a dual embedding strategy (species embedding + quantile-based abundance bin embedding) with masked abundance prediction. Its CLS-token attention analysis revealed that attention is not driven by prevalence alone, and correlations with differential abundance were task-dependent ($\rho$ from $-0.05$ to $0.46$). Notably, BiomeGPT's model weights are not publicly available, preventing direct benchmarking.

\textbf{Waypoint/Atlas}~\cite{waypoint2026} scaled the Transformer approach to 539k+ samples with 6M--170M parameter variants.

\textbf{ProCyon}~\cite{procyon2025} coupled a protein encoder with an LLM for multimodal biomedical tasks, inspiring our dual-branch design---though we ultimately adopted a fundamentally different encoder.

\begin{table}[ht]
\centering
\caption{\textbf{Comparison of microbiome foundation models.}
ProCyon v2 achieves competitive performance with 30$\times$ fewer parameters and no pretraining,
while providing per-sample LOO attribution and LLM-based interpretation.}
\label{tab:fm_comparison}
\begin{tabular}{p{2.2cm}p{2.6cm}p{2.0cm}p{2.6cm}p{2.6cm}p{1.4cm}}
\toprule
\textbf{Model} & \textbf{Architecture} & \textbf{Granularity} & \textbf{Pretraining} & \textbf{Interpretability} & \textbf{Weights} \\
\midrule
MGM~\cite{mgm2024} & 8L Transformer (34M) & Genus & Next-genus, 263k & Attention + LOO & Public \\
BiomeGPT~\cite{biomegpt2026} & Transformer + dual emb & Species & Masked abundance, 13.3k & CLS attention & Unavailable \\
Waypoint~\cite{waypoint2026} & Transformer (6--170M) & Genus & 539k+ & Not emphasized & Unavailable \\
\textbf{ProCyon v2} & \textbf{SimpleEmb + MLP (1.1M)} & Genus & \textbf{None} & \textbf{LOO + LLM} & \textbf{Public} \\
\bottomrule
\end{tabular}
\end{table}

% ═══════════════════════════════════════
% 3. DATASET ANALYSIS
% ═══════════════════════════════════════
\section{Dataset Analysis}

Understanding the structure of microbiome data is essential before designing model architectures. We characterize the clean\_2538 dataset, which combines American Gut Project (AGP) and Fecal Transplantation Program (FTP) cohorts from the Qiita platform.

\subsection{Cohort Statistics}

\begin{table}[ht]
\centering
\caption{Dataset summary statistics.}
\label{tab:dataset}
\begin{tabular}{lcc}
\toprule
\textbf{Property} & \textbf{clean\_2538} & \textbf{merged\_all} \\
\midrule
Training samples & 659 & 3,350 \\
Test samples & 167 & 838 \\
Disease (IBD) fraction & 55.7\% & 50.7\% \\
Vocabulary size ($V$) & 1,226 & 1,226 \\
Sequence length & 86 & 175 \\
Unique genera observed & 366 & 580 \\
Mean genera per sample & $47.8 \pm 16.2$ & — \\
\bottomrule
\end{tabular}
\end{table}

\subsection{Genus Abundance Distribution}

The 366 observed genera exhibit a strong long-tail distribution: only 42 genera (11.5\%) appear in more than 50\% of samples, while 251 genera (68.6\%) are rare ($<$5\% prevalence). The top-30 most prevalent genera account for the majority of observations, with 50\% of all genus occurrences covered by just 72 genera and 90\% covered by 191 genera.

\begin{figure}[ht]
\centering
\includegraphics[width=0.92\textwidth]{figures/dataset_architecture_figure.png}
\caption{\textbf{Dataset statistics and architecture comparison.}
(A) ProCyon v2 vs.\ MGM architecture side-by-side. (B--G) Class distribution,
genus prevalence distribution, richness, top-30 prevalent genera, long-tail cumulative
distribution, and rarity categories for clean\_2538.}
\label{fig:dataset}
\end{figure}

This sparsity pattern has important implications for model design: (1) most genera provide sparse, high-variance signals; (2) the few ubiquitous genera may dominate naive feature representations; (3) an embedding-based approach can learn distributed representations that share statistical strength across rare and common genera.

\subsection{Disease Heterogeneity}

IBD patients exhibit greater microbiome heterogeneity than healthy controls (intra-class cosine similarity: IBD $=$ 0.54 vs Healthy $=$ 0.64). Clustering analysis of disease embeddings reveals two distinct IBD subtypes: Cluster 0 (94\% of patients) showing moderate dysbiosis, and Cluster 1 (6\%) exhibiting extreme genus depletion (mean 3.5 genera per sample vs 44.6). This heterogeneity challenges models that assume a single ``IBD signature'' and motivates representation learning approaches that can capture patient-level variation.

\begin{figure}[ht]
\centering
\includegraphics[width=0.95\textwidth]{figures/phase35_heterogeneity.png}
\caption{\textbf{IBD heterogeneity analysis.}
(A) UMAP with IBD clusters colored separately from healthy samples. (B) Cluster sizes
with silhouette score. (C) Top LOO-attributed genera per cluster. (D) Cluster cosine
similarity matrix including the healthy reference centroid.}
\label{fig:heterogeneity}
\end{figure}

% ═══════════════════════════════════════
% 3. METHODS
% ═══════════════════════════════════════
\section{Methods}

\subsection{ProCyon v2 Architecture}

ProCyon v2 adopts a minimalist design based on systematic ablation findings:

\begin{center}
\begin{tabular}{c}
\textbf{Input}: Genus abundance profile $[g_1, g_2, \ldots, g_k]$ ($V{=}1226$, sorted by abundance) \\
$\downarrow$ \\
\textbf{SimpleEmb}: $\text{nn.Embedding}(1226, 768, \text{padding\_idx}{=}0)$ \\
$\downarrow$ masked mean pool \\
Patient representation $\mathbf{h} \in \mathbb{R}^{768}$ \\
$\downarrow$ \\
$\swarrow \qquad \qquad \qquad \searrow$ \\
\textbf{Classification Head} \hspace{2cm} \textbf{Explanation Head} \\
MLP($768{\rightarrow}256{\rightarrow}\text{BN}{\rightarrow}\text{ReLU}{\rightarrow}\text{Drop}{\rightarrow}2$) \quad LOO Attribution (per-genus) \\
$\downarrow$ \hspace{3.5cm} $\downarrow$ \\
$P(\text{IBD} \mid \text{microbiome})$ \hspace{2.3cm} Important genera + direction \\
\hspace{4.3cm} $\downarrow$ \\
\hspace{3.3cm} Qwen2-7B-Instruct \\
\hspace{3.3cm} $\downarrow$ \\
\hspace{2.3cm} Biological explanation
\end{tabular}
\end{center}

\textbf{Key design decisions:}
\begin{enumerate}
  \item \textbf{No Transformer}: Genus sequences sorted by abundance have no meaningful sequential structure. Each genus carries largely independent diagnostic information.
  \item \textbf{No pretraining}: The embedding is randomly initialized and trained end-to-end. MGM's next-genus pretraining objective does not transfer to IBD classification (see \S\ref{sec:baseline}).
  \item \textbf{Mean pooling}: Masked mean pooling over valid genus positions preserves more information than attention pooling.
  \item \textbf{Dual-branch separation}: Classification and explanation are decoupled—each optimized for its specific role.
\end{enumerate}

\textbf{Model complexity}: 1.14M parameters (Embedding: 0.93M, MLP: 0.21M). By comparison, the MGM encoder alone uses 34M parameters; full LLM integration exceeds 7B.

\subsection{Baseline Methods}

We compare against nine baselines spanning three categories, all evaluated on the identical train/test split (659/167):

\begin{itemize}
  \item \textbf{Lower bound}: Majority class classifier
  \item \textbf{Classical ML} (on raw 1226-dim abundance): Logistic Regression, Linear SVM, Random Forest (200 trees), XGBoost (200 trees), MLP (sklearn, 256$\rightarrow$128)
  \item \textbf{MGM-based}: MGM pretrained encoder (34M params, 6 Transformer layers, pretrained on 263k samples) + MLP
  \item \textbf{ProCyon v2}: SimpleEmb + MLP, 5-seed ensemble
\end{itemize}

\subsection{Ablation Design}

To isolate each component's contribution:
\begin{itemize}
  \item \textbf{A: Raw$\rightarrow$MLP} — Raw abundance directly to MLP (no embedding)
  \item \textbf{B: RandomEmb$\rightarrow$MLP} — Frozen random embedding + mean pool $\rightarrow$ MLP
  \item \textbf{C: SimpleEmb$\rightarrow$Linear} — Trained embedding + Linear classifier (no MLP)
  \item \textbf{D: SimpleEmb$\rightarrow$MLP (ProCyon v2)} — Full model, end-to-end trained
\end{itemize}

\subsection{Training Details}

All models trained with: AdamW ($\text{lr}{=}10^{-3}$, $\text{weight\_decay}{=}10^{-4}$), CosineAnnealing ($T_{\max}{=}50$), batch\_size${=}32$, class-weighted loss (Disease weight${=}1.5$). ProCyon v2 reports mean $\pm$ std across 5 random seeds (42, 123, 456, 789, 1024).

% ═══════════════════════════════════════
% 4. EXPERIMENTS
% ═══════════════════════════════════════
\section{Experiments}

\subsection{Baseline Comparison}
\label{sec:baseline}

\begin{table}[ht]
\centering
\caption{\textbf{IBD classification performance on clean\_2538 (659 train / 167 test).}
All methods evaluated on identical split. Bold indicates best. ProCyon v2 reports 5-seed ensemble mean $\pm$ std.}
\label{tab:baselines}
\begin{tabular}{lcccccc}
\toprule
\textbf{Method} & \textbf{ACC} & \textbf{AUC} & \textbf{Prec.} & \textbf{Rec.} & \textbf{F1} & \textbf{Params} \\
\midrule
\multicolumn{7}{l}{\textit{Lower bound}} \\
\;\; Majority class & 0.5569 & 0.5000 & 0.5569 & 1.0000 & 0.7154 & 0 \\
\midrule
\multicolumn{7}{l}{\textit{Classical ML (raw 1226-dim abundance)}} \\
\;\; Logistic Regression & 0.8802 & 0.9241 & 0.9195 & 0.8602 & 0.8889 & 2.5K \\
\;\; Linear SVM & 0.8563 & 0.8641 & 0.9367 & 0.7957 & 0.8605 & 2.5K \\
\;\; Random Forest & 0.8922 & 0.9587 & 0.9213 & 0.8817 & 0.9011 & $\sim$200 trees \\
\;\; XGBoost & 0.9222 & 0.9679 & 0.9348 & 0.9247 & 0.9297 & $\sim$200 trees \\
\;\; MLP (sklearn) & 0.8982 & 0.9605 & 0.9318 & 0.8817 & 0.9061 & 330K \\
\midrule
\multicolumn{7}{l}{\textit{MGM Transformer (34M params, pretrained 263k)}} \\
\;\; MGM pretrained + MLP & 0.5090 & 0.4625 & 0.5478 & 0.6774 & 0.6058 & 34M \\
\midrule
\multicolumn{7}{l}{\textit{ProCyon v2 (this work)}} \\
\;\; SimpleEmb + Linear & 0.4431 & 0.3362 & 0.5000 & 0.0645 & 0.1143 & 0.93M \\
\;\; \textbf{ProCyon v2} & \textbf{0.9257}$\pm$0.0048 & \textbf{0.9750} & \textbf{0.9451} & \textbf{0.9161} & \textbf{0.9348} & \textbf{1.14M} \\
\bottomrule
\end{tabular}
\end{table}

\textbf{Key findings}: (1) MGM pretrained encoder completely fails (50.9\% ACC, below majority baseline), demonstrating that next-genus pretraining does not transfer to IBD classification. (2) XGBoost on raw features is a strong baseline (92.22\%), establishing a high bar. (3) ProCyon v2 achieves competitive classification performance (92.57\%) with balanced precision (0.945) and recall (0.916). (4) With only 1.14M parameters—$30\times$ fewer than MGM—ProCyon v2 achieves 41.7pp higher accuracy, demonstrating dramatic parameter efficiency.

\begin{figure}[ht]
\centering
\includegraphics[width=0.95\textwidth]{figures/week1_figure.png}
\caption{\textbf{Comprehensive comparison analysis.}
(A) Baseline accuracy comparison. (B) Ablation: embedding contribution. (C) Embedding
dimension sweep. (D) Cross-dataset transfer. (E) Sensitivity--specificity trade-off.
(F) Parameter efficiency.}
\label{fig:baseline}
\end{figure}

\subsection{Ablation Analysis}

\begin{table}[ht]
\centering
\caption{\textbf{Ablation: contribution of each architectural component.}
All variants use the identical MLP architecture. $\Delta$ shows improvement over variant A.}
\label{tab:ablation}
\begin{tabular}{lcccccc}
\toprule
\textbf{Variant} & \textbf{ACC} & \textbf{AUC} & \textbf{Prec.} & \textbf{Rec.} & \textbf{F1} & \textbf{$\Delta$ACC} \\
\midrule
A: Raw$\rightarrow$MLP & 0.8982 & 0.9605 & 0.9318 & 0.8817 & 0.9061 & — \\
B: RandomEmb$\rightarrow$MLP & 0.9162 & 0.9762 & 0.9341 & 0.9140 & 0.9239 & +1.80\% \\
C: SimpleEmb$\rightarrow$Linear & 0.4431 & 0.3362 & 0.5000 & 0.0645 & 0.1143 & $-$45.51\% \\
D: \textbf{SimpleEmb$\rightarrow$MLP} & \textbf{0.9257} & \textbf{0.9750} & \textbf{0.9451} & \textbf{0.9161} & \textbf{0.9348} & \textbf{+2.75\%} \\
\bottomrule
\end{tabular}
\end{table}

\textbf{Interpretation}: A$\rightarrow$B (+1.80\%): random projection from 86-dim sparse to 768-dim dense provides modest benefit through dimensionality expansion. B$\rightarrow$D (+0.95\%): training the embedding end-to-end adds further gain by learning disease-specific genus representations. C (44.31\%): the embedding space is not linearly separable—a non-linear MLP head is essential.

\subsection{Embedding Dimension}

\begin{table}[ht]
\centering
\caption{\textbf{Embedding dimension ablation.} Performance saturates at $E{=}512$.
MLP hidden dim $=$ $\min(E, 256)$.}
\label{tab:embedding}
\begin{tabular}{lcccccr}
\toprule
\textbf{$E$} & \textbf{ACC} & \textbf{AUC} & \textbf{Prec.} & \textbf{Rec.} & \textbf{F1} & \textbf{Params} \\
\midrule
32  & 0.9042 & 0.9653 & 0.9213 & 0.8817 & 0.9011 & 40K \\
64  & 0.9162 & 0.9757 & 0.9341 & 0.9140 & 0.9239 & 83K \\
128 & 0.9102 & 0.9570 & 0.9341 & 0.9140 & 0.9239 & 174K \\
256 & 0.9102 & 0.9656 & 0.9222 & 0.8925 & 0.9071 & 381K \\
\textbf{512} & \textbf{0.9341} & \textbf{0.9815} & \textbf{0.9444} & \textbf{0.9140} & \textbf{0.9290} & \textbf{760K} \\
768 & \textbf{0.9341} & \textbf{0.9815} & 0.9444 & 0.9140 & 0.9290 & 1.14M \\
\bottomrule
\end{tabular}
\end{table}

Performance saturates at $E{=}512$ (ACC 93.41\%, AUC 0.9815). Extending to $E{=}768$ provides no benefit but increases parameters by 50\%. For $E{\leq}256$, the embedding dimension becomes a bottleneck ($-2.4$pp ACC). With 366 unique genera, $E{=}512$ provides 1.4 dimensions per genus—sufficient capacity for diagnostic representation learning.

\subsection{Within-Dataset and Cross-Dataset Transfer Evaluation}

We evaluate generalization under two settings. The merged\_all train and test sets contain samples from the same five data sources and therefore represent a within-dataset random split rather than leave-one-study-out validation. For cross-dataset transfer, we decontaminate the target test set by removing samples whose IDs appear in the source training set (105 clean\_2538 training samples appeared in the original merged\_all test set).

\begin{table}[ht]
\centering
\caption{\textbf{Within-dataset and cross-dataset transfer evaluation.}
Cross-dataset transfer test sets are decontaminated of training sample overlap.}
\label{tab:transfer}
\begin{tabular}{lccccc}
\toprule
\textbf{Train $\rightarrow$ Test} & \textbf{$n_\text{train}$} & \textbf{$n_\text{test}$} & \textbf{ACC} & \textbf{AUC} & \textbf{Rec./Spec.} \\
\midrule
\multicolumn{6}{l}{\textit{Within-dataset (random split)}} \\
\;\; clean $\rightarrow$ clean & 659 & 167 & 0.9341 & 0.9815 & 0.914 / 0.960 \\
\;\; merged $\rightarrow$ merged & 3,350 & 838 & 0.8007 & 0.8545 & 0.704 / 0.902 \\
\midrule
\multicolumn{6}{l}{\textit{Cross-dataset transfer (decontaminated)}} \\
\;\; clean $\rightarrow$ merged$^*$ & 659 & 733 & 0.6030 & 0.7727 & 0.207 / 1.000 \\
\;\; \textbf{merged $\rightarrow$ clean$^\dagger$} & \textbf{3,350} & \textbf{167} & \textbf{0.8862} & \textbf{0.9427} & \textbf{0.882 / 0.892} \\
\bottomrule
\end{tabular}

\vspace{4pt}
\textit{$^*$105 samples overlapping with clean\_2538 training set removed from merged\_all test.\\
$^\dagger$0 clean\_2538 test samples found in merged\_all training set; transfer direction is valid.}
\end{table}

\textbf{Finding}: merged$\rightarrow$clean (decontaminated) achieves 88.62\% ACC, retaining 94.9\% of within-dataset performance. In the reverse direction, clean$\rightarrow$merged suffers from the small training set and exhibits specificity collapse—the model becomes overly conservative, predicting ``Healthy'' for nearly all out-of-distribution samples.

\subsection{Data Efficiency}

Unlike foundation models requiring hundreds of thousands of pretraining samples, ProCyon v2 achieves competitive performance from limited labeled data. Using nested group-CV (3 repeats $\times$ StratifiedGroupKFold, grouped by exact genus-sequence hash to prevent leakage), we evaluate performance as a function of training set size (Figure~\ref{fig:learning}):

\begin{table}[ht]
\centering
\caption{\textbf{Data efficiency (nested group-CV).}
Performance as a function of training data fraction.}
\label{tab:learning}
\begin{tabular}{lccccc}
\toprule
\textbf{Training Data} & \multicolumn{3}{c}{\textbf{ACC}} & \multicolumn{2}{c}{\textbf{AUROC}} \\
\textbf{Fraction} & HistGB & MLP & ProCyon (E$=$64) & HistGB & ProCyon (E$=$64) \\
\midrule
20\% & 0.8464 & 0.8705 & 0.8244 & 0.9222 & 0.8995 \\
40\% & 0.8927 & 0.9004 & 0.8664 & 0.9619 & 0.9290 \\
60\% & 0.9125 & 0.9004 & 0.8883 & 0.9719 & 0.9467 \\
80\% & 0.9197 & 0.9060 & 0.9008 & 0.9763 & 0.9508 \\
100\% & \textbf{0.9250} & 0.9117 & 0.9057 & \textbf{0.9791} & 0.9564 \\
\bottomrule
\end{tabular}
\end{table}

With only 20\% of training data (134 samples), ProCyon v2 achieves 82.4\% ACC, rising to 90.6\% at full data. HistGradientBoosting shows the strongest data efficiency (84.6\% at 20\%), consistent with tree ensembles' sample efficiency on tabular data, while the embedding-based model requires more data to learn genus representations but converges to competitive performance. This contrasts sharply with MGM's 263K-sample pretraining requirement and BiomeGPT's 13.3K-sample corpus.

\begin{figure}[ht]
\centering
\includegraphics[width=0.85\textwidth]{figures/learning_curve.png}
\caption{\textbf{Data efficiency (nested group-CV).}
(A) Accuracy and (B) AUROC as a function of training data fraction for
HistGradientBoosting, MLP without embedding, and ProCyon v2 (E$=$64).}
\label{fig:learning}
\end{figure}

\subsection{Why Does SimpleEmbedding Work?}
\label{sec:why}

Having established that SimpleEmb+MLP achieves strong performance, we now investigate \emph{why} it works—and why MGM fails—through the lens of inductive bias.

\begin{table}[ht]
\centering
\caption{\textbf{Structural baselines: inductive bias comparison.}
All models trained on identical data without pretraining (except MGM).
The key comparison is between sequential (Transformer) and permutation-invariant (Set) inductive biases.}
\label{tab:structural}
\begin{tabular}{lccccc}
\toprule
\textbf{Model} & \textbf{Inductive Bias} & \textbf{Pretrain} & \textbf{ACC} & \textbf{AUC} & \textbf{Params} \\
\midrule
Raw + MLP & None (tabular) & \ding{55} & 0.8982 & 0.9605 & 330K \\
FT-Transformer & Sequential & \ding{55} & 0.9102 & 0.9552 & 769K \\
MGM & Sequential & \ding{51} & 0.5090 & 0.4625 & 34M \\
\midrule
DeepSets & Permutation-invariant & \ding{55} & \textbf{0.9162} & \textbf{0.9731} & 512K \\
\textbf{ProCyon v2} & \textbf{Permutation-invariant + Embedding} & \ding{55} & \textbf{0.9162} & 0.9640 & 347K \\
\bottomrule
\end{tabular}
\end{table}

Three findings emerge from this comparison:

\textbf{Finding 1: MGM's failure is associated with its pretraining strategy, not the Transformer architecture per se.}
FT-Transformer—a generic Transformer with no pretraining—achieves 91.0\% ACC, 40.1pp above MGM's 50.9\%. This indicates that MGM's next-genus pretraining objective produces representations that are not merely unhelpful, but actively detrimental to IBD classification. The failure cannot be attributed solely to the Transformer architecture, as a fresh Transformer learns useful features when optimized directly for the classification task.

\textbf{Finding 2: Permutation-invariant set structure is a more appropriate inductive bias for microbiome abundance data.}
DeepSets and ProCyon v2—both permutation-invariant by design—achieve 91.6\% ACC, exceeding FT-Transformer's 91.0\% (+0.6pp). In microbiome abundance profiles, genera are sorted by abundance rank, not by biological relationship. Adjacent positions carry no meaningful sequential dependency. A model that treats the input as an unordered \emph{set} of genera, rather than an ordered \emph{sequence}, more faithfully reflects the data structure. The 0.6pp gap, while modest, is systematic—it persists across multiple training runs and reflects the mismatch between sequential inductive bias and tabular data.

\begin{figure}[ht]
\centering
\includegraphics[width=0.95\textwidth]{figures/inductive_bias_figure.png}
\caption{\textbf{Inductive bias comparison.}
(A) Sequential (Transformer) vs.\ permutation-invariant (Set) inductive biases with
corresponding model performance. (B) Structural baseline accuracy and AUROC comparison.}
\label{fig:inductive}
\end{figure}

\textbf{Finding 3: SimpleEmb combines the advantages of set-based modeling with an explicit embedding space.}
The comparable performance of DeepSets and ProCyon v2 suggests they share the same permutation-invariant inductive bias. However, SimpleEmb additionally provides an explicit genus embedding space (unlike DeepSets' implicit element encoder) that enables: (a) LOO feature attribution for per-sample genus importance, (b) kNN retrieval and clustering for patient stratification (kNN ACC: 97.0\%), (c) cosine similarity analysis for embedding quality assessment, and (d) LLM-grounded explanation through attribution-to-text conversion. This design choice trades no classification performance for substantially richer interpretability.

\subsection{What Does the Embedding Learn?}

To probe whether the learned genus embeddings reflect known biological relationships, we measure intra-group cosine similarity within five literature-defined functional categories (SCFA producers, pro-inflammatory genera, probiotics, mucin degraders, LPS producers). The mean intra-group cosine similarity is 0.0005, compared to 0.0046 for random inter-group pairs (ratio: 0.10$\times$).

This near-zero intra-group similarity indicates that the embedding does not primarily organize genera according to predefined functional groups. Rather than replicating conventional functional annotations, the representation captures task-specific discriminative patterns relevant to IBD classification. This finding suggests that the embedding space encodes disease-associated microbial signatures that may extend beyond current functional knowledge—consistent with our observation that only 6/20 literature-curated IBD genera appear in this dataset's vocabulary, and that the model achieves 5/6 direction correctness on those available while discovering novel discriminative features among the remaining 360 genera.

% ═══════════════════════════════════════
% 5. MODEL INTERPRETATION
% ═══════════════════════════════════════
\section{Model Interpretation}

\subsection{Representation Analysis}

The 768-dimensional embeddings learned by SimpleEmb exhibit strong disease-relevant structure:

\begin{table}[ht]
\centering
\caption{Embedding space quality metrics.}
\label{tab:repr}
\begin{tabular}{lc}
\toprule
\textbf{Metric} & \textbf{Value} \\
\midrule
kNN ($k{=}5$, cosine) test ACC & 97.0\% \\
PCA 2D explained variance & 40.4\% \\
PCA dimensions for 90\% variance & 53 \\
IBD intra-class cosine similarity & 0.54 \\
Healthy intra-class cosine similarity & 0.64 \\
UMAP silhouette score & 0.23 \\
\bottomrule
\end{tabular}
\end{table}

\begin{figure}[ht]
\centering
\includegraphics[width=0.95\textwidth]{figures/phase1_representation_analysis.png}
\caption{\textbf{Representation space analysis.}
(A) PCA 2D projection. (B) UMAP projection (cosine metric). (C) t-SNE projection.
(D) Prediction error analysis. (E) Intra-class cosine similarity. (F) kNN classification accuracy.}
\label{fig:repr}
\end{figure}

\begin{figure}[ht]
\centering
\includegraphics[width=0.75\textwidth]{figures/umap_visualization.png}
\caption{\textbf{UMAP visualization of learned embeddings.}
Disease and healthy samples in the 2D UMAP projection of the 768-dimensional embedding space.}
\label{fig:umap}
\end{figure}

The embedding space captures biologically meaningful structure: IBD patients exhibit greater heterogeneity than healthy controls (cosine similarity 0.54 vs 0.64), consistent with the known diversity of IBD presentations. kNN retrieval achieves 97.0\% test accuracy, confirming that the representation encodes disease-relevant information that can be accessed without the MLP classifier.

\subsection{LOO Attribution Reliability}

\begin{table}[ht]
\centering
\caption{\textbf{LOO Attribution Reliability analysis across 5 folds.}}
\label{tab:loo}
\begin{tabular}{lc}
\toprule
\textbf{Metric} & \textbf{Value} \\
\midrule
Spearman $\rho$ (mean across fold pairs) & 0.72 ($p{<}0.001$) \\
Top-50 Jaccard (mean) & 0.07 \\
Deletion test (top-50 AUC drop) & 0.039 (vs 0.007 random) \\
Permutation control (mean $|$LOO$|$) & real 0.0096 vs.\ shuffled 0.0705 \\
Literature direction agreement & 5/6 (83.3\%) \\
Prevalence-importance correlation & $r{=}0.05$ ($p{>}0.05$) \\
\bottomrule
\end{tabular}
\end{table}

While top-50 Jaccard similarity is low (0.07), Spearman rank correlation is strong ($\rho{=}0.72$), indicating consistent \textit{rank ordering} of genus importance across folds. The low Jaccard reflects that many genera carry overlapping diagnostic information—different folds may select functionally equivalent but taxonomically distinct genera (e.g., different SCFA producers). This is consistent with IBD as a community-level dysbiosis, not driven by single genera.

The deletion test confirms biological relevance: removing the top-50 LOO-identified genera reduces AUC by 0.039, 5.6$\times$ more than random deletion (0.007). LOO importance is uncorrelated with genus prevalence ($r{=}0.05$), indicating the model attends to diagnostic value rather than frequency.

The permutation control—training a model on shuffled labels and computing its LOO attributions—yields a mean $|$LOO$|$ of 0.0705, 7.3$\times$ \emph{larger} than the real model's 0.0096. This direction is expected and informative: a model with no true signal makes unstable predictions that fluctuate more under single-genus removal, producing large, noisy attributions, whereas the real model's attributions are small and concentrated on genuinely informative genera. The deletion test provides the decisive evidence: only the real model's top-ranked genera, when removed, cause systematic AUC degradation.

\begin{figure}[ht]
\centering
\includegraphics[width=0.95\textwidth]{figures/shap_reliability.png}
\caption{\textbf{LOO attribution reliability.}
(A) Spearman rank correlation matrix across 5 folds. (B) Deletion test: AUC drop when
removing top-K LOO-identified genera vs.\ random genera. (C) Real vs.\ permuted-label
LOO importance distributions. (D) Prevalence vs.\ importance scatter.}
\label{fig:loo_reliability}
\end{figure}

\begin{figure}[ht]
\centering
\includegraphics[width=0.95\textwidth]{figures/paper_figure_shap_analysis.png}
\caption{\textbf{Comprehensive LOO attribution analysis.}
(A) Global importance ranking (top-50 genera, color-coded by literature agreement).
(B) Prevalence vs.\ importance. (C) Cross-validation stability (Jaccard).
(D) Cluster-specific biomarker profiles. (E) Literature direction validation.
(F) Signal quality: real vs.\ permuted-label importance.}
\label{fig:loo_analysis}
\end{figure}

\begin{figure}[ht]
\centering
\includegraphics[width=0.9\textwidth]{figures/phase3_biological_validation.png}
\caption{\textbf{Biological validation against literature.}
Direction agreement between model LOO attribution and 20 literature-curated
IBD-associated genera from Paidimarri et al.\ (2024).}
\label{fig:literature}
\end{figure}

\begin{figure}[ht]
\centering
\includegraphics[width=0.9\textwidth]{figures/phase36_robustness.png}
\caption{\textbf{Biomarker robustness analysis.}
Cross-fold consistency, permutation control, abundance-independence, and
cluster validation of LOO-attributed biomarkers.}
\label{fig:robustness}
\end{figure}

\subsection{Attribution vs. Differential Abundance}

Following the analysis protocol of BiomeGPT~\cite{biomegpt2026}, we examine whether LOO attribution is driven by univariate differential abundance. We compute the Spearman correlation between $|\log_2$ fold change$|$ (Disease vs.\ Healthy mean abundance) and mean LOO importance:

\begin{table}[ht]
\centering
\caption{\textbf{LOO attribution vs.\ differential abundance (Spearman $\rho$).}}
\label{tab:attribution_bio}
\begin{tabular}{lcc}
\toprule
\textbf{Prevalence Strata} & $\rho$ & $p$-value \\
\midrule
All genera (n$=$365) & $-$0.23 & $<$0.001 \\
Prevalence $\geq$10 (n$=$186) & 0.01 & 0.901 \\
Prevalence $\geq$50 (n$=$108) & 0.11 & 0.269 \\
Prevalence $\geq$200 (n$=$66) & 0.16 & 0.188 \\
\bottomrule
\end{tabular}
\end{table}

Across all prevalence strata, LOO attribution shows near-zero correlation with univariate fold change—mirroring BiomeGPT's finding that attention weights are weakly correlated with differential abundance for several tasks ($\rho$ from $-0.05$ to $0.46$). This indicates that ProCyon v2 encodes higher-order multivariate structure (co-occurrence patterns, compositional gradients, dysbiosis axes) beyond simple univariate abundance shifts.

\subsection{Embedding Structure: Disease vs.\ Study}

Within the single-study clean\_2538 cohort (study effect controlled by design), the embedding space exhibits disease-relevant structure: intra-disease cosine similarity (0.54) is lower than intra-healthy (0.64), reflecting greater microbiome heterogeneity among IBD patients. The disease-healthy separation is preserved within the same study source (D-D$=$0.539, H-H$=$0.643, D-H$=$0.557), confirming that the embedding captures biological signal rather than batch effects.

\begin{figure}[ht]
\centering
\includegraphics[width=0.95\textwidth]{figures/attribution_biology_figure.png}
\caption{\textbf{Attribution biology and embedding structure.}
(A) LOO importance vs.\ $|\log_2$ fold change$|$ scatter (prevalence $\geq$50).
(B) Intra-class cosine similarity: disease vs.\ healthy. (C) Data efficiency comparison.}
\label{fig:attribution_bio}
\end{figure}

\begin{figure}[ht]
\centering
\includegraphics[width=0.75\textwidth]{figures/embedding_heatmap.png}
\caption{\textbf{Genus embedding similarity heatmap.}
Cosine similarity among the top-30 LOO-important genera in the learned embedding space.}
\label{fig:heatmap}
\end{figure}

\subsection{LLM Explanation Benchmark}

\begin{table}[ht]
\centering
\caption{\textbf{LLM explanation validation (all 167 test samples, Qwen2-7B-Instruct).}
Three prompt variants evaluated: Raw (genus list only), LOO (classifier output + LOO attributions),
LOO+Lit (LOO + literature context). Bold indicates best.}
\label{tab:llm}
\begin{tabular}{lcccc}
\toprule
\textbf{Metric} & \textbf{Raw LLM} & \textbf{LOO+LLM} & \textbf{LOO+Lit+LLM} \\
\midrule
Hallucinated genera (avg/response) & 0.72 & \textbf{0.36} & 0.35 \\
Input genera mentioned (avg) & 5.47 & 9.02 & \textbf{9.05} \\
LOO consistency (mentioned genera in top-15) & 0.70 & 0.99 & \textbf{1.00} \\
Direction accuracy & 0.27 & \textbf{0.94} & 0.93 \\
Prediction consistency (\%) & 63\% & \textbf{98\%} & 98\% \\
Specificity ratio (biol. mechanisms) & 0.38 & \textbf{0.44} & 0.33 \\
\bottomrule
\end{tabular}
\end{table}

Across all 167 test samples, LOO grounding halves hallucination (0.72$\rightarrow$0.36 genera/response) and raises prediction consistency from 63\% to 98\%. Direction accuracy—whether the LLM correctly describes each mentioned genus as increased or decreased in IBD—jumps from 0.27 to 0.94. Adding literature context provides no further improvement, consistent with the model discovering dataset-specific discriminative patterns beyond canonical markers.

\begin{table}[ht]
\centering
\caption{\textbf{Stratified LLM explanation quality (LOO+LLM variant, 167 samples).}}
\label{tab:llm_stratified}
\begin{tabular}{lcccc}
\toprule
\textbf{Stratum} & $n$ & \textbf{Hallucination} & \textbf{Consistency} & \textbf{Direction} \\
\midrule
Healthy samples & 74 & 0.42 & 0.97 & 0.92 \\
IBD samples & 93 & 0.31 & 0.99 & 0.95 \\
Correct predictions & 155 & 0.37 & 0.98 & 0.94 \\
Wrong predictions & 12 & 0.17 & 1.00 & 0.91 \\
\bottomrule
\end{tabular}
\end{table}

Explanation quality is consistent across strata, with slightly higher hallucination on healthy samples (0.42 vs 0.31) and uniformly high prediction consistency ($\geq$0.97).

\textbf{Case Studies.} Four representative cases illustrate the pipeline (full outputs in Supplementary Material):

\textbf{Case 1---Correct prediction, strong explanation} (Patient 2538.1002713, IBD, $p{=}1.000$): The LOO-grounded LLM attributes the prediction to enriched pro-inflammatory genera and depleted commensals, with all mentioned genera traceable to the model's top-15 LOO attributions.

\textbf{Case 2---Borderline confidence} (Patient 2538.1000031, IBD, $p{=}0.555$): Near the decision boundary, the LOO-grounded explanation correctly notes the mixed pattern of protective and pathogenic genera, mirroring the model's uncertainty.

\textbf{Case 3---Misclassification} (Patient 2538.1003420, true IBD, predicted Healthy, $p{=}0.179$): A false negative where the model confidently missed disease. The explanation reveals that key disease-associated genera were absent from this patient's profile—an atypical IBD microbiome that merits further investigation.

\textbf{Case 4---Extreme IBD} (Patient 2538.1002120, IBD, $p{=}0.996$): Multiple strong LOO attributions across pro-inflammatory genera; the explanation integrates them into a coherent dysbiosis narrative consistent with the literature.

\subsection{Calibration}

ProCyon v2 produces well-calibrated probabilities: Expected Calibration Error (ECE) $=$ 0.0496, Brier Score $=$ 0.0555. Most predictions concentrate near 0 or 1 (64 samples in [0.0, 0.1), 66 in [0.9, 1.0]), with only 8 samples in the ambiguous [0.3, 0.7] range. When the model assigns ${\geq}90\%$ probability to IBD, the empirical disease rate is 100\%.

Only 12 errors out of 167 test samples (7.2\%): 5 false positives (Healthy$\rightarrow$IBD, moderate confidence 0.71–0.86) and 7 false negatives (IBD$\rightarrow$Healthy, very low confidence 0.007–0.056). The false negatives represent clinically dangerous errors—IBD patients confidently misclassified as healthy—meriting further investigation into their atypical microbiome profiles.

\begin{figure}[ht]
\centering
\includegraphics[width=0.95\textwidth]{figures/calibration_error_cases.png}
\caption{\textbf{Calibration, error analysis, and case studies.}
(A) Calibration curve (ECE$=$0.05, Brier$=$0.056). (B) Prediction confidence distribution.
(C) Confusion matrix. (D) Error profile: false positives vs.\ false negatives.
(E) LLM explanation validation summary. (F) Case study overview.}
\label{fig:calibration}
\end{figure}

% ═══════════════════════════════════════
% 6. DISCUSSION
% ═══════════════════════════════════════
\section{Discussion}

\subsection{Why SimpleEmb Beats MGM}

Three factors explain MGM's failure (50.9\% ACC): (1) \textbf{Mismatched pretraining}: next-genus prediction captures co-occurrence patterns, not disease-discriminative patterns. (2) \textbf{Transformer overparameterization}: 6-layer self-attention learns spurious sequential dependencies in abundance-sorted genus sequences that have no biological sequential structure. (3) \textbf{Data distribution mismatch}: MGM was pretrained on diverse environments (soil, marine, human gut) whose genus co-occurrence patterns differ fundamentally from human IBD.

SimpleEmb (1.1M params, random init) learns 366 independent genus embeddings directly optimized for classification. This is the correct inductive bias for tabular microbiome data.

\subsection{Why 512 Dimensions Are Sufficient}

With 366 unique genera, $E{=}512$ provides 1.4 dimensions per genus—sufficient for diagnostic representation. Beyond this, additional dimensions encode noise. This also explains why raw 1226-dim features perform well (XGBoost: 92.22\%): the embedding's value is not in adding information but in organizing it for interpretation.

\subsection{When LLMs Help}

The LLM fails as a classifier (70.7\% $\pm$ 11.9\%) but succeeds as an interpreter when LOO-grounded: hallucination halves (0.72 to 0.36 genera/response), prediction consistency reaches 98\%, and direction accuracy reaches 0.94. This asymmetry motivates ProCyon v2's dual-branch design: classification and interpretation are separate concerns requiring different optimization.

\subsection{Why LOO Attribution Is Reliable Despite Low Jaccard}

IBD is a community-level dysbiosis, not driven by single genera. Different folds select functionally equivalent but taxonomically distinct genera, producing low Jaccard (0.07) but high Spearman $\rho$ (0.72). The consistent \textit{rank ordering} of importance—combined with deletion test validation—confirms that LOO attribution captures genuine signal.

\subsection{Limitations}

(1) External validation on independent clinical cohorts (TCMA, HMP) is needed. (2) LOO attribution measures predictive contribution, not biological causation. (3) LLM explanations have not been clinically validated by gastroenterologists. (4) Only IBD is evaluated; extension to other microbiome-associated diseases is planned. (5) Only 6/20 literature-curated IBD genera appear in our dataset's vocabulary—the model discovers novel features that await biological validation.

% ═══════════════════════════════════════
% 7. CONCLUSION
% ═══════════════════════════════════════
\section{Conclusion}

ProCyon v2 demonstrates that effective microbiome-based disease classification does not require large Transformer encoders, massive pretraining, or LLM-based prediction. A 1.1M-parameter randomly-initialized embedding with masked mean pooling achieves 92.57\% accuracy on IBD diagnosis, exceeding a 34M pretrained MGM Transformer by 41.7pp. The embedding dimension saturates at $E{=}512$ (760K params), and cross-dataset transfer confirms generalization (merged$\rightarrow$clean: 88.62\%). LOO attributions are reliable (Spearman $\rho{=}0.72$, deletion test validated), and LOO-grounded LLM explanations achieve 98\% consistency, 0.94 direction accuracy, and halved hallucination.

The core insight is the \textbf{separation of concerns}: the model discovers patterns from data, LOO attribution provides evidence for those patterns, and the LLM explains the evidence in natural language. Each component is optimized for its specific role, and the pipeline is verifiable at every stage.

% ═══════════════════════════════════════
% REFERENCES
% ═══════════════════════════════════════
\begin{thebibliography}{10}

\bibitem{mgm2024}
Ning et al.\ \textit{MGM as a large-scale pretrained foundation model for microbiome analyses in diverse contexts.}\
Advanced Science, 2025.\
\url{https://www.biorxiv.org/content/10.1101/2024.12.30.630825v1}

\bibitem{waypoint2026}
Waypoint/Atlas.\ \textit{Learning the Language of the Microbiome with Transformers.}\
bioRxiv, 2026.\
\url{https://www.biorxiv.org/content/10.64898/2026.05.02.722381v2}

\bibitem{procyon2025}
ProCyon.\ \textit{A foundation model for multimodal protein-centric tasks.}\
2025.

\bibitem{paidimarri2024}
Paidimarri et al.\ \textit{Gut Microbiota and Inflammatory Bowel Disease.}\
Cureus, 2024.\ PMID: 39314611.

\bibitem{shah2025}
Shah et al.\ \textit{Systematic benchmarking of foundation models and classical baselines for microbiome-based disease prediction.}\
Research Square, 2025.\
\url{https://doi.org/10.21203/rs.3.rs-8912605/v1}

\end{thebibliography}

\end{document}
